\chapter{F\_pulsalyzer\_real\_time\_visualizer.py}
\label{ch:scriptF}

This script sends rotation commands to Pulsalyzer using the \texttt{A\_json\_commander.py} script, thus it can be used to visualize the effect of rotation on the Pulsalyzer-rendered point cloud in real-time.

\section{How It Works}
Pulsalyzer accepts a scanner transformation command (\texttt{set\_scanner\_transformation\_deg}) via its JSON WebServer. This script sends the transformation command along with \texttt{reprocess\_pointcloud} and \texttt{redraw\_pointcloud} commands through \texttt{A\_json\_commander.py}, incrementally updating rotation angle along one axis and triggering a reprocess and redraw of the point cloud in Pulsalyzer.

The commands are sent through a \texttt{send\_silent()} function which suppresses terminal output from \texttt{A\_json\_commander.py}, keeping the terminal clean during the loop.

\section{User-Configurable Parameters}
\begin{table}[!ht]
    \centering
    \begin{tabularx}{\textwidth}{llX}
        \toprule
        Variable               & Default & Description                                                                \\
        \midrule
        routing1               & [...]   & Routing path to the Pulsalyzer endpoint; must include the correct hostname \\
        x\_rot, y\_rot, z\_rot & 0, 0, 0 & Initial rotation angles (degrees)                                          \\
        start, stop            & 0, 10   & Start and end values of the angle increment range                          \\
        \bottomrule
    \end{tabularx}
\end{table}

\section{Prerequisites}
\begin{itemize}
    \item JSONServer and Pulsalyzer must be running, with Pulsalyzer having the required \texttt{.lidar} file and its corresponding trajectory loaded and processed.
    \item \texttt{A\_json\_commander.py} script must be in the same directory as this script.
    \item The hostname must match the device name visible in the Pulsalyzer WebServer interface (or \texttt{hostname} command in command prompt).
\end{itemize}

\section{Usage}

The script applies the initial rotation angles, pauses for a keypress, then enters an infinite loop cycling through the angle increment range. The loop restarts when the range is exhausted. Press \texttt{Ctrl+C} to stop the loop.

By default, the increment is applied to the Z-axis rotation, modify as needed for other axes. The update rate of point cloud in the 3D viewer of Pulsalyzer depends on the \texttt{point\_step} parameter of commands \texttt{reprocess\_pointcloud} and \texttt{redraw\_pointcloud}.

This script was used in the thesis to identify the hidden rotation angle applied internally by Pulsalyzer. By iterating through a range of rotation angles (\qty{-180}{\degree} to \qty{180}{\degree}) about the X and Z axes and observing the rendered point cloud in real-time, the angles at which the Pulsalyzer output visually matched the Python translation-only reference output was identified as \qty{180}{\degree}, \qty{0}{\degree}, \qty{90}{\degree} for X, Y, and Z axes respectively. This hidden rotation matrix was then used to derive the mathematical relationship between the Python and Pulsalyzer boresight angles.